\documentclass[aspectratio=169]{beamer}

% ---------- THEME & LOOK ----------
\usetheme{Madrid}
\usecolortheme{seahorse}
\usefonttheme{professionalfonts}
\setbeamertemplate{navigation symbols}{}
\setbeamercovered{transparent}
\setbeamertemplate{itemize items}[circle]
\setbeamertemplate{enumerate items}[default]
\setlength{\parskip}{0.25em}

% ---------- COLORS ----------
\usepackage{xcolor}
\definecolor{accent}{RGB}{0,92,184}
\definecolor{soft}{RGB}{233,244,255}
\setbeamercolor{structure}{fg=accent}
\setbeamercolor{title}{fg=white,bg=accent}
\setbeamercolor{frametitle}{fg=accent}
\setbeamercolor{block title}{bg=accent!12,fg=accent}
\setbeamercolor{block body}{bg=soft}

% ---------- PACKAGES ----------
\usepackage{booktabs}
\usepackage{amsmath, amssymb}
\usepackage{array}
\usepackage{hyperref}
\hypersetup{colorlinks=true,linkcolor=accent,urlcolor=accent}

% ---------- PLACEHOLDER BOX (compile-safe; no images required) ----------
\newcommand{\placeholderbox}[2][0.82\linewidth]{%
  \begin{center}
    \fbox{\begin{minipage}[c][0.46\textheight][c]{#1}
      \centering \vspace{0.25em}
      \textbf{#2}\\[0.6em]
      \emph{Insert figure/diagram here later.}
    \end{minipage}}
  \end{center}
}

% ---------- TITLE ----------
\title{\textbf{Human Capital Theory: Applications to Education and Training \vspace{.5cm}\\  }}
\subtitle{ {\large ECON 300 - Labour Economics}}
%\institute{UNBC}
\author{Muhebullah Karimzada}

\date{Winter 2026}


\begin{document}

% ============================================================
% TITLE
% ============================================================
\begin{frame}[plain]
  \titlepage
\end{frame}

% ============================================================
\section{Orientation}
% ============================================================

\begin{frame}{Learning Objectives}
\transfade
\begin{itemize}[<+->]
  \item Analyze schooling and on-the-job training as standard investments: compare costs and benefits.
  \item Recognize the \textbf{opportunity cost of time} as the largest cost.
  \item Understand why simple earnings comparisons can mislead (ability bias; signalling/screening).
  \item Summarize how returns to education in Canada are large and have risen over time.
  \item Identify when workers versus firms should pay for training.
\end{itemize}
\end{frame}

\begin{frame}{Main Questions (1/2)}
\transfade
\begin{itemize}[<+->]
  \item Why are more-educated workers generally paid more?
  \item What determines the market rate of return to education?
  \item If education pays, why doesn’t everyone pursue the maximum degree?
  \item Human capital vs.\ signalling/screening: which story fits the data?
\end{itemize}
\end{frame}

\begin{frame}{Main Questions (2/2)}
\transfade
\begin{itemize}[<+->]
  \item How might education serve as a signal and why might private/social returns diverge?
  \item Should government subsidize early childhood learning programs?
  \item Are government-funded training programs worth the money?
\end{itemize}
\end{frame}

% ============================================================
\section{Human Capital Theory}
% ============================================================

\begin{frame}{Human Capital Theory: Core Ideas}
\transfade
\begin{itemize}[<+->]
  \item Investments made to \textbf{raise productivity and earnings}.
  \item Costs today (tuition, materials, \textbf{foregone earnings}); benefits tomorrow (higher income).
  \item Undertake the investment when \textbf{benefits exceed costs}.
\end{itemize}
\end{frame}

\begin{frame}{Costs and Benefits: Types}
\transfade
\begin{itemize}[<+->]
  \item \textbf{Opportunity cost}: value of time in alternative employment.
  \item \textbf{Private vs.\ Social} costs/benefits.
  \item \textbf{Real costs} vs.\ \textbf{transfer costs}.
\end{itemize}
\end{frame}

% ============================================================
\section{Private Investment in Education}
% ============================================================

\begin{frame}{Figure 9.1: Education and Alternative Income Streams — \textit{Placeholder}}
\transfade
\placeholderbox{Figure 9.1 — Age–earnings profiles with vs.\ without further schooling.}
\end{frame}

\begin{frame}{Age–Earnings Profiles}
\transfade
\begin{itemize}[<+->]
  \item Increase with age but at a \textbf{decreasing rate}.
  \item Higher for those with \textbf{more education}.
\end{itemize}
\end{frame}

\begin{frame}{Optimal Investment: Assumptions}
\transfade
\begin{itemize}[<+->]
  \item No direct utility/disutility from education (consumption value ignored).
  \item Hours of work are fixed.
  \item Income streams by education level are known.
  \item Perfect capital markets: can borrow/lend at real interest rate.
\end{itemize}
\end{frame}

\begin{frame}{NPV of Benefits: Setup}
\transfade
\small
Consider an 18-year-old HS graduate choosing work vs.\ college.\par
Let working life be \(T-18\) years and constant income \(Y\).\par
\textbf{Present value of the income stream at 18:}
\[
\text{PV} = \sum_{t=1}^{T-18} \frac{Y}{(1+r)^t}
\]
where \(r\) is the market (discount) rate.
\end{frame}

\begin{frame}{Decision Rules: Margins and IRR}
\transfade
\small
\textbf{Cost of one more year (MC):}
\[
MC = Y + D
\]
\textbf{Benefit of one more year (MB):}
\[
MB = \sum_{t=1}^{T-19}\frac{\Delta Y}{(1+r)^t}
\]
Optimal years of schooling when \(\textbf{MB}=\textbf{MC}\).\par
\textbf{Internal rate of return (IRR):} choose schooling until \(i=r\).
\end{frame}

\begin{frame}{Figure 9.2: Two Ways to State the Decision Rule — \textit{Placeholder}}
\transfade
\placeholderbox{Figure 9.2 — (1) MB=MC in PV terms; (2) Internal rate of return i equals market r.}
\end{frame}

\begin{frame}{Implications and Influences}
\transfade
\begin{itemize}[<+->]
  \item Invest \textbf{early} to enjoy more years of higher earnings.
  \item Less incentive with \textbf{interrupted careers}.
  \item Progressive taxes can reduce demand for education.
  \item Student loans alter marginal costs and attainment.
\end{itemize}
\end{frame}

% ============================================================
\section{Worked Example: Returns to Education}
% ============================================================

\begin{frame}{Worked Example: Computing Returns}
\transfade
\small
Mary (18) can work now at \$30{,}000 or take \textbf{one year} of post-secondary.\par
Interest rate \(r=5\%\). Work until 65.\par
After school, income becomes \$33{,}000.\par
\textbf{Approximation:} For a level annuity, \( \text{PV} \approx \frac{Y}{r} \) (infinite horizon approximation).\par
The extra year raises PV of income by roughly \$23{,}943.\par
\textbf{IRR:} \(i = \frac{\Delta Y}{Y + D}\). If \(D=0\), \(i=3000/30000=10\%\). If \(D=3000\), \(i\approx 9\%\).\par
Direct costs would need to rise to about \$30{,}000 for Mary not to invest.
\end{frame}

% ============================================================
\section{Education as Signal / Filter}
% ============================================================

\begin{frame}{Education as a Filter}
\transfade
\begin{itemize}[<+->]
  \item Education \textbf{signals} productivity to employers.
  \item Higher wages offered if education raises productivity \emph{or} credibly reveals higher ability.
\end{itemize}
\end{frame}

\begin{frame}{Figure 9.3: Offered Wage and Signalling Cost Schedules — \textit{Placeholder}}
\transfade
\placeholderbox{Figure 9.3 — Offered wage vs.\ signalling cost; separating equilibrium intuition.}
\end{frame}

% ============================================================
\section{Empirical Evidence: Education and Earnings}
% ============================================================

\begin{frame}{Stylized Facts}
\transfade
\begin{itemize}[<+->]
  \item Earnings rise with experience; most rapid to ages ~40–44 for most-educated.
  \item Differentials widen with age (wider at 50 than 20–30).
\end{itemize}
\end{frame}

\begin{frame}{Figure 9.4: Earnings by Age and Education, Canadian Males, 2015 — \textit{Placeholder}}
\transfade
\placeholderbox{Figure 9.4 — Age–earnings profiles by education group, males (2015).}
\end{frame}

\begin{frame}{Table 9.1: Estimates of the Private Returns\textsuperscript{1} to Schooling in Canada, 2000}
\transfade
\scriptsize
\begin{tabular}{lcc}
\toprule
\textbf{Level of Schooling} & \textbf{Males} & \textbf{Females} \\
\midrule
Bachelor’s degree\textsuperscript{2} & 12 & 14 \\
Master’s degree & 3 & 5 \\
Ph.D. & nc\textsuperscript{2} & 4 \\
Medicine & 21 & 22 \\
\midrule
\textbf{Bachelor’s Degree by Field of Study} & \textbf{Males} & \textbf{Females} \\
Education & 9 & 14 \\
Humanities and fine arts & nc & 10 \\
Social sciences\textsuperscript{3} & 11 & 14 \\
Commerce & 9 & 19 \\
Natural sciences & 9 & 8 \\
Engineering and applied science & 9 & 14 \\
Health sciences & 18 & 18 \\
\bottomrule
\end{tabular}

\vspace{0.25em}
\scriptsize \textit{Notes:} \textsuperscript{1} Private monetary returns. \textsuperscript{2} “nc” reported where not computed. \textsuperscript{3} Includes economics and related fields.
\end{frame}

\begin{frame}{Human Capital Earnings Function}
\transfade
\small
\[
\ln Y = \alpha + rS + \beta_1 \text{EXP} + \beta_2 \text{EXP}^2 + \varepsilon
\]
\begin{itemize}[<+->]
  \item \(Y\): earnings, \(S\): years of schooling (return \(r\)), \(\text{EXP}\): experience proxy, \(\varepsilon\): unobservables.
  \item Controls help separate schooling effects from experience and ability.
\end{itemize}
\end{frame}

\begin{frame}{Figure 9.5: Log Earnings by Years of Schooling (Women 35–39, 2015) — \textit{Placeholder}}
\transfade
\placeholderbox{Figure 9.5 — Log earnings vs.\ years of schooling with fitted line.}
\end{frame}

\begin{frame}{Table 9.2: Estimated Returns to Schooling and Experience, 2015\\(Dependent Variable: Log Annual Earnings)}
\transfade
\scriptsize
\begin{tabular}{lcc}
\toprule
 & \textbf{Men} & \textbf{Women} \\
\midrule
Intercept & \;9.232 \;(714.15) & \;8.605 \;(628.29) \\
Years of schooling & \;0.080 \;(105.31) & 0.114 \;(140.77) \\
Experience & 0.054 \;(75.96) & 0.041 \;(62.51) \\
Experience squared & $-0.0009$ \;(60.14) & $-0.0006$ \;(43.85) \\
$R^2$ & 0.136 & 0.204 \\
Sample size & 121{,}947 & 99{,}398 \\
\bottomrule
\end{tabular}
\end{frame}

\begin{frame}{Ability Bias, Signalling, and Screening}
\transfade
\begin{itemize}[<+->]
  \item Difficult to fully control for innate ability, motivation, perseverance.
  \item Natural experiments used to address bias: twins; compulsory schooling laws; proximity to college.
  \item Evidence suggests earlier fears of large ability bias were \textbf{overstated}.
\end{itemize}
\end{frame}

\begin{frame}{Private vs.\ Social Returns}
\transfade
\begin{itemize}[<+->]
  \item Private estimates focus on monetary benefits to the individual.
  \item Social returns may differ (spillovers: innovation, civic outcomes, crime reduction).
\end{itemize}
\end{frame}

\begin{frame}{Increased Returns and Inequality}
\transfade
\begin{itemize}[<+->]
  \item Returns rose substantially since the mid-1990s.
  \item Polarization: strong gains at the very top.
  \item Inequality reflects both returns across and \emph{within} education groups.
\end{itemize}
\end{frame}

% ============================================================
\section{Training}
% ============================================================

\begin{frame}{Training: General vs.\ Specific}
\transfade
\begin{itemize}[<+->]
  \item \textbf{General training:} portable skills; trainees accept lower current wages to recoup later.
  \item \textbf{Specific training:} firm-specific; costs and benefits often shared since mobility is limited.
\end{itemize}
\end{frame}

\begin{frame}{Figure 9.6: Costs, Benefits, and Financing of Training — \textit{Placeholder}}
\transfade
\placeholderbox{Figure 9.6 — Profiles of wages/costs under general vs.\ specific training.}
\end{frame}

\begin{frame}{Training and the Role of Government}
\transfade
\begin{itemize}[<+->]
  \item Markets may underprovide training due to imperfect information or constraints.
  \item Targeted subsidies can raise hours and lift wages above poverty thresholds.
\end{itemize}
\end{frame}

% ============================================================
\section{Chapter Summary}
% ============================================================

\begin{frame}{Summary}
\transfade
\begin{itemize}[<+->]
  \item Human capital investment: compare PV benefits with costs; stop when \(MB=MC\) or \(i=r\).
  \item Education shifts age–earnings profiles; earnings functions quantify returns.
  \item Signalling/screening complicate naïve comparisons; empirical strategies address ability bias.
  \item On-the-job training: who pays depends on portability (general vs.\ specific).
  \item Policy: early education, access to finance, and targeted training can improve outcomes.
\end{itemize}
\end{frame}

% ============================================================
\begin{frame}
\centering
\Huge{\textbf{Thank You!}} \\

\end{frame}

\end{document}
